\documentclass[10pt,a4paper]{amsart}
\usepackage[margin=19mm]{geometry}
\usepackage{amsmath,amssymb,mathtools}
\usepackage[scr=rsfs,cal=euler]{mathalfa}
\usepackage{xfrac}
\usepackage[unicode,hidelinks,bookmarks=false]{hyperref}
\newcommand{\ie}{i.e.\ }
\newcommand{\cf}{cf.\ }
\newcommand{\resp}{respectively\ }
\newcommand{\Subsection}{\subsection}
\providecommand{\nobreakdash}{\nobreak\hbox{-}}
\setlength{\emergencystretch}{2em}
\multlinegap0pt
\usepackage{bm}
\usepackage{bbm}
\usepackage{dsfont}
\usepackage{xspace}
\usepackage{cancel}
\usepackage{mathtools}
\usepackage{amsthm}
\makeatletter
\@ifundefined{theorem}{\newtheorem{theorem}{Theorem}}{}
\@ifundefined{lemma}{\newtheorem{lemma}[theorem]{Lemma}}{}
\makeatother
\newcommand{\B}{\mathcal{B}}
\newcommand{\E}{\mathbb{E}}
\newcommand{\F}{\mathcal{F}}
\newcommand{\N}{\mathbb{N}}
\newcommand{\X}{\mathcal{X}}
\newcommand{\Y}{\mathcal{Y}}
\newcommand{\dint}{\mathrm{d}}  % integrálás d-betű
\newcommand{\eps}{\varepsilon}
\newcommand{\ga}{\gamma}
\newcommand{\ind}{\mathds{1}}
\newcommand{\ka}{\kappa}
\newcommand{\leb}{\mathrm{Leb}}
\newcommand{\lfrf}[1]{\left\lfloor #1\right\rfloor}
\renewcommand{\P}{\mathbb{P}}
\title[Transition of $\alpha$-mixing in Random Iterations with Appl]{Transition of $\alpha$-mixing in Random Iterations with Applications in Queuing Theory}
\author{Attila Lovas}
\date{Source: arXiv:2410.05056v5 (CC BY 4.0)}
\begin{document}
\maketitle

\noindent\emph{Source and licence.} Transition of $\alpha$-mixing in Random Iterations with Applications in Queuing Theory. Version arXiv:2410.05056v5 (CC BY 4.0), distributed under the Creative Commons Attribution 4.0 International licence, as recorded in the arXiv licence metadata for this exact version and shown on the abstract page https://arxiv.org/abs/2410.05056v5. Full rights notice: licenses/arxiv-2410.05056-CC-BY-4.0.txt.\par\smallskip

\noindent\textit{Editorial scope.} Counted unit: the Lemma whose source label is \texttt{lem:coupling} in the article (source0.tex lines 2072--2085, source label \texttt{lem:coupling}), delivered together with the article's own complete original proof of it (source0.tex lines 2086--2257, one \texttt{proof} environment of 172 lines). No mathematics is written, repaired or re-derived: statement and proof are the article's own text line for line. Required context retained uncounted: eq:Lyapunov (lines 382--384); eq:smallset (lines 390--392); lem:T (lines 1966--1981); eq:fgood (lines 2057--2060); eq:Zaux (lines 2064--2069). Every definition, display and label of those lines is retained unchanged. Omitted from this package and not redistributed: 1--381, 385--389, 393--1965, 1982--2056, 2061--2063, 2070--2071, 2258--2499. Nothing inside a retained excerpt is omitted or altered: every retained body is delivered exactly as the article's lines are written, so no omission receipt of any kind accompanies this package. Citations: the retained excerpts carry no citation command at all, so no bibliography entry of the article is transcribed and no reference is redistributed. Reference adaptation: none was needed and none was made; every reference inside the delivered unit resolves inside it, so no unresolved reference or citation is produced. Printed slips kept literal: no printed word, symbol, hypothesis or number of the counted statement or of its proof was altered. Layout and class: the article's own class is \texttt{article} with packages including inputenc, authblk, amsthm, amsmath, amssymb, latexsym, dsfont, enumerate, tikz, xcolor, url, geometry; the wrapper is the lane's pinned \texttt{amsart} editorial wrapper at 10pt on A4, which restates the theorem-like environments the retained text needs and adds the article's own macro definitions that the retained text needs (\texttt{\textbackslash B}, \texttt{\textbackslash E}, \texttt{\textbackslash F}, \texttt{\textbackslash N}, \texttt{\textbackslash P}, \texttt{\textbackslash X}, \texttt{\textbackslash Y}, \texttt{\textbackslash dint}, \texttt{\textbackslash eps}, \texttt{\textbackslash ga}, \texttt{\textbackslash ind}, \texttt{\textbackslash ka}, \texttt{\textbackslash leb}, \texttt{\textbackslash lfrf}), with extra wrapper packages (bm, bbm, dsfont, xspace, cancel, mathtools). The delivered numbering is the unit's own. Assets: no retained span contains a figure environment, a graphics-inclusion command, a PDF-page inclusion, a table or a TikZ picture; no image file is copied into this package, the package declares no asset dependency and no image file is redistributed with it. Licence: version arXiv:2410.05056v5 of this article is distributed under the Creative Commons Attribution 4.0 International licence, as recorded in the arXiv licence metadata for this exact version and shown on the abstract page https://arxiv.org/abs/2410.05056v5. A full rights notice is delivered at licenses/arxiv-2410.05056-CC-BY-4.0.txt. Characters: every retained excerpt is pure ASCII, so the delivered engine, XeLaTeX with its default Latin Modern text font, reports no missing character.
\begin{equation}\label{eq:Lyapunov}
[Q(y)V](x):=\int_{\X}V(z)\,Q(y,x,\dint z)\le \ga (y)V(x)+K(y).
\end{equation}

\begin{equation}\label{eq:smallset}
Q(y,x,A)\ge (1-\beta (R,y))\ka_R (y,A). 
\end{equation}

\begin{lemma}\label{lem:T}
	Let $R>0$ be arbitrary and suppose the parametric kernel $Q:\Y\times\X\times\B(\X)\to [0,1]$ satisfies the minorization condition given by \eqref{eq:smallset}. Then there exists a measurable mapping $T^R:\X\times\Y\times [0,1]\to\X$ such that
	$$
	Q(y,x,A)=\int_{[0,1]}\ind_{T^R (x,y,u)\in A}\,\dint u,
	$$
	for all $x\in\X$, $A\in\B(\X)$ and $y\in\Y$. 
	Furthermore, for any fixed $y\in\Y$, there exists a Borel set $U=U (y)\in \B ([0,1])$
	with Lebesgue measure
	$
	\leb_1 (U)\ge 1-\beta (R,y)
	$
	such that for $u\in U$,
	\begin{equation}\label{eq:Tconst}
	T^R (x_1,y,u) = T^R (x_2,y,u),\,\,x_1,x_2\in\stackrel{-1}{V}([0,R]).
	\end{equation}
\end{lemma}

\begin{equation}\label{eq:fgood}
	(x,y,u)\mapsto f(x,y,u):=T^{R(y)}(x,y,u),\,\,
	x\in\X,\,y\in\Y,\,u\in [0,1].
\end{equation}

\begin{equation}\label{eq:Zaux}
Z_{s,t}^{x,\mathbf{y}} = \begin{cases}
x &\text{ if } t\le s \\
f(Z_{s,t-1}^{x,\mathbf{y}},y_{t-1},\eps_t) &\text{ if } t>s,
\end{cases}
\end{equation}

\begin{lemma}\label{lem:coupling}
	Let $x_1,x_2\in\X$ and $\mathbf{y}\in\Y^\N$ be arbitrary but fixed. Then for $0<m<n$, we have
	$$
	\P (Z_{0,n}^{x_1,\mathbf{y}} \ne Z_{0,n}^{x_2,\mathbf{y}})
	\le
	\bar{\beta}^m
	+
	(1+r)^{\lfrf{\frac{n-1}{m}}}\left\{
	\frac{V(x_1)+V(x_2)}{2}\prod_{j=0}^{\lfrf{\frac{n-1}{m}}}\ga (y_j)
	+
	\sum_{k=0}^{m-1}\sum_{l=k}^{k\lfrf{\frac{n-1}{m}}} K(y_l)\prod_{j=1}^{\lfrf{\frac{n-1}{m}}}\ga (y_{l+j})
	\right\}.
	$$
\end{lemma}

\begin{proof}
	For the sake of brevity, we employ a more concise notation: $Z_n^i=Z_{0,n}^{x_i,\mathbf{y}}$, $i=1,2$, and $n\in\N$, moreover we introduce
	$$
	\overline{Z}_{n}:=\left(Z_n^1,Z_n^2\right), \,\,
	\lVert \overline{Z}_{n}\rVert:=\max \left( V(Z_n^1), V(Z_n^2) \right)
	$$ 
	and the sequence of successive visiting times
	$$
	\sigma_0:=0,\,\sigma_{k+1} = \inf\left\{n>\sigma_k\middle|  \lVert \overline{Z}_n\rVert\le 
	R(y_n)
	\right\},\,k\in\N
	$$
	that are obviously $\F_{1,\infty}^\eps$-stopping times. Note that on 
	$\{ \lVert \overline{Z}_n\rVert>R (y_n)\}$ we have 
	\begin{equation}\label{eq:contr}
	\ga (y_n) (V( Z_n^1)+V( Z_n^2))+2K(y_n)
	\le
	(1+r)\ga (y_n) (V( Z_n^1)+V( Z_n^2))
	\end{equation}
	and thus for $k\ge 1$ and $s\ge 1$, by the Markov inequality, we obtain
	\begin{align*}
	\P (\sigma_{k+1}-\sigma_{k}>s\mid \F_{1,\sigma_k}^\eps)
	&\le 
	\E \left(
	\P (C_{\sigma_k+s}\mid \overline{Z}_{\sigma_{k}+s-1})
	\prod_{j=1}^{s-1}
	\ind_{C_{\sigma_k+j}}
	\middle| \F_{1,\sigma_k}^\eps\right)\\
	&\le 
	\frac{(1+r)\ga (y_{\sigma_k+s-1})}{R(y_{\sigma_k+s})}
	\E \left(
	\left(V( Z_{\sigma_k+s-1}^1)+V( Z_{\sigma_k+s-1}^2)\right)
	\prod_{j=1}^{s-2}
	\ind_{C_{\sigma_k+j}}
	\middle| \F_{1,\sigma_k}^\eps\right),
	\end{align*}
	where $C_n$ is the shorthand notation for the event $\{\lVert \overline{Z}_n\rVert> 
	R(y_n)\}$, $n\in\N$. 
	
	By the tower rule, we can write
	\begin{multline*}
		\E \left(
		\left(V( Z_{\sigma_k+s-1}^1)+V( Z_{\sigma_k+s-1}^2)\right)
		\prod_{j=1}^{s-2}
		\ind_{C_{\sigma_k+j}}
		\middle| \F_{1,\sigma_k}^\eps\right)
		=\\
		\E \left(
		\E\left[V( Z_{\sigma_k+s-1}^1)+V( Z_{\sigma_k+s-1}^2)\middle|\F_{1,\sigma_{k}+s-2}^\eps\right]
		\prod_{j=1}^{s-2}
		\ind_{C_{\sigma_k+j}}
		\middle| \F_{1,\sigma_k}^\eps\right).
	\end{multline*}
	
	Using the Markov property of $(\bar{Z}_n)_{n\in\N}$ and the drift property of the parametric kernel \eqref{eq:Lyapunov}, we have
	\begin{align*}
		\E\left[V( Z_{\sigma_k+s-1}^1)+V( Z_{\sigma_k+s-1}^2)\middle|\F_{1,\sigma_{k}+s-2}^\eps\right]
		&=
		\E\left[V( Z_{\sigma_k+s-1}^1)+V( Z_{\sigma_k+s-1}^2)\middle|\bar{Z}_{\sigma_{k}+s-2}\right]
		\\
		&=
		[Q(y_{\sigma_k+s-2})V](Z_{\sigma_k+s-2}^1)+
		[Q(y_{\sigma_k+s-2})V](Z_{\sigma_k+s-2}^2)
		\\
		&\le
		\ga (y_{\sigma_k+s-2})\left(V(Z_{\sigma_k+s-2}^1) + (Z_{\sigma_k+s-2}^2)\right)+2K(y_{\sigma_k+s-2}).
	\end{align*}
	Now, inequality \eqref{eq:contr} yields
	$$
	\E\left[V( Z_{\sigma_k+s-1}^1)+V( Z_{\sigma_k+s-1}^2)\middle|\F_{1,\sigma_{k}+s-2}^\eps\right]\ind_{C_{\sigma_k+s-2}}\le (1+r)\ga (y_{\sigma_k+s-2})\left(V(Z_{\sigma_k+s-2}^1) + (Z_{\sigma_k+s-2}^2)\right).
	$$
	Finally, we arrive at
	\begin{multline*}
			\E \left(
		\left(V( Z_{\sigma_k+s-1}^1)+V( Z_{\sigma_k+s-1}^2)\right)
		\prod_{j=1}^{s-2}
		\ind_{C_{\sigma_k+j}}
		\middle| \F_{1,\sigma_k}^\eps\right)
		\le 
		\\
		(1+r)\ga (y_{\sigma_k+s-2})
		\E \left(
		\left(V( Z_{\sigma_k+s-2}^1)+V( Z_{\sigma_k+s-2}^2)\right)
		\prod_{j=1}^{s-3}
		\ind_{C_{\sigma_k+j}}
		\middle| \F_{1,\sigma_k}^\eps\right).
	\end{multline*}
	
	Iteration of this argument in $s-2$ steps leads to the following estimation:
	\begin{align*}
	\P (\sigma_{k+1}-\sigma_{k}>s\mid \F_{1,\sigma_k}^\eps) 
	&\le  
	\frac{(1+r)^{s-1}\prod_{j=1}^{s-1}\ga (y_{\sigma_k+j})}{R(y_{\sigma_k+s})}
	\E \left(
	V( Z_{\sigma_k+1}^1)+V( Z_{\sigma_k+1}^2)
	\middle| \overline{Z}_{\sigma_k}\right) \\
	&\le
	\frac{r (1+r)^{s-1}\prod_{j=1}^{s}\ga (y_{\sigma_k+j})}{2K(y_{\sigma_k+s})}
	\left[\ga (y_{\sigma_k})\left(V( Z_{\sigma_k}^1)+V( Z_{\sigma_k}^2)\right) +2K(y_{\sigma_k})\right] \\
	&\le (1+r)^s K(y_{\sigma_k})\prod_{j=1}^{s}\ga (y_{\sigma_k+j}),
	\end{align*}
	where we used that $V( Z_{\sigma_k}^1)+V( Z_{\sigma_k}^2)\le 2R(y_{\sigma_k})$, and $K(\cdot)\ge 1$.
	
	Along similar lines, we can show that 
	\begin{align*}
		\P (\sigma_1>s) &\le \frac{(1+r)^{s-1}\prod_{j=1}^{s-1}\ga (y_j)}{R(y_s)}
		\left[\ga (y_0)\left(V(x_1)+V(x_2)\right)+2K(y_0)\right]
		\\
		&\le
		(1+r)^s \prod_{j=1}^{s}\ga (y_j)\left[\frac{\ga (y_0)}{2}\left(V(x_1)+V(x_2)\right)+K(y_0)\right].
	\end{align*}
	
	Clearly, for any $0<m\le n$, on the event $\{\sigma_m>n\}$ we have $\{\sigma_{k+1}-\sigma_{k}>\lfrf{n/m}\}\cap\{\sigma_{k}\le k\lfrf{n/m}\}$ for some $k=0,1,\ldots m-1$ hence by the union bound and the estimates we obtain for the time elapsed between consecutive visits, we can write
	\begin{align*}
		\P (\sigma_m>n)&\le \P \left(\bigcup_{k=0}^{m-1} \{\sigma_{k+1}-\sigma_{k}>\lfrf{n/m}\}\cap\{\sigma_{k}\le k\lfrf{n/m}\}\right)
		\\
		&\le 
		\sum_{k=0}^{m-1}\P \left(\{\sigma_{k+1}-\sigma_{k}>\lfrf{n/m}\}\cap\{\sigma_{k}\le k\lfrf{n/m}\}\right)
		\\
		&=
		\sum_{k=0}^{m-1}\sum_{l=k}^{k\lfrf{n/m}}
		\P \left(\sigma_{k+1}-\sigma_{k}>\lfrf{n/m}\mid \sigma_{k}=l\right)\P(\sigma_{k}=l)
		\\
		&\le 
		\P(\sigma_{1}>\lfrf{n/m})+
			\sum_{k=1}^{m-1}\sum_{l=k}^{k\lfrf{n/m}}
		\P \left(\sigma_{k+1}-\sigma_{k}>\lfrf{n/m}\mid \sigma_{k}=l\right)
	\\
	&\le 
	\left[\frac{\ga (y_0)}{2}\left(V(x_1)+V(x_2)\right)+K(y_0)\right]
	(1+r)^{\lfrf{n/m}} \prod_{j=1}^{\lfrf{n/m}}\ga (y_j)
	\\
	&+
	(1+r)^{\lfrf{n/m}}\sum_{k=1}^{m-1}\sum_{l=k}^{k\lfrf{n/m}} K(y_l)\prod_{j=1}^{\lfrf{n/m}}\ga (y_{l+j}).
	\end{align*}
	
	Next, we estimate the probability of no-coupling on events when small sets are visited at least $m$-times. By Lemma \ref{lem:T} and the choice of $R(y)$ in the definition of $f$ in \eqref{eq:fgood} and \eqref{eq:Zaux}, for each $j=1,\ldots,m$, $x \mapsto f(x,y_j,\eps_{\sigma_j+1})$ is constant on the level set $\stackrel{-1}{V}([0,R(y_{\sigma_j})])$ with probability at least $1-\bar{\beta}$ 
	independently of $\F_{0,\sigma_j}^\eps$ thus no-coupling happens with probability at most $\bar{\beta}$.
	Therefore, we can estimate:
	\begin{align*}
	\P (Z_{0,n}^{x_1,\mathbf{y}}\ne Z_{0,n}^{x_2,\mathbf{y}},\sigma_{m}<n) 
	\le \bar{\beta}^m. 
	\end{align*}
	
	Finally, we combine this upper bound with that one what we got for the tail probability of the visiting times, and obtain
	\begin{align*}
	&\P (Z_{0,n}^{x_1,\mathbf{y}}\ne Z_{0,n}^{x_2,\mathbf{y}})
	\le 
	\P (Z_{0,n}^{x_1,\mathbf{y}}\ne Z_{0,n}^{x_2,\mathbf{y}},\sigma_{m}<n)
	+
	\P (\sigma_{m}> n-1)\\
	& \le
	\bar{\beta}^m
	+
	\left[\frac{\ga (y_0)}{2}\left(V(x_1)+V(x_2)\right)+K(y_0)\right]
	(1+r)^{\lfrf{\frac{n-1}{m}}} \prod_{j=1}^{\lfrf{\frac{n-1}{m}}}\ga (y_j)
	\\
	&+
	(1+r)^{\lfrf{\frac{n-1}{m}}}\sum_{k=1}^{m-1}\sum_{l=k}^{k\lfrf{\frac{n-1}{m}}} K(y_l)\prod_{j=1}^{\lfrf{\frac{n-1}{m}}}\ga (y_{l+j})
	\\
	&=
	\bar{\beta}^m
	+
	(1+r)^{\lfrf{\frac{n-1}{m}}}\left\{
	\frac{V(x_1)+V(x_2)}{2}\prod_{j=0}^{\lfrf{\frac{n-1}{m}}}\ga (y_j)
	+
	\sum_{k=0}^{m-1}\sum_{l=k}^{k\lfrf{\frac{n-1}{m}}} K(y_l)\prod_{j=1}^{\lfrf{\frac{n-1}{m}}}\ga (y_{l+j})
	\right\}
	\end{align*}
	which completes the proof.
	
\end{proof}

\end{document}
